\documentclass[../main.tex]{subfiles}

\begin{document}

\chapter{$L^{p}$ 空间: 对偶与弱收敛}\label{chapter:ch08-对偶与弱收敛}

对于可测集 $E$, $1\leqslant p<\infty$, 而 $q$ 为 $p$ 的共轭, 令 $g$ 属于 $L^{q}(E)$. 定义 $L^{p}(E)$ 上的实值泛函 $T$ 为
\begin{align*}
\text{对所有 }f\in L^{p}(E),\quad T(f)=\int_{E}f\cdot g.
\end{align*}
H\"older 不等式告诉我们 $f\cdot g$ 是可积的, 因此 $T$ 是正确定义的. 泛函 $T$ 继承了积分的线性. 此外, 存在常数 $M\geqslant 0$ 使得对所有 $f\in L^{p}(E)$,
\begin{align*}
|T(f)|\leqslant M\cdot\|f\|_{p}.
\end{align*}
事实上, 根据 H\"older 不等式, 这对 $M=\|g\|_{q}$ 成立. Riesz 表示定理断言, 若 $T$ 是 $L^{p}(E)$ 上的任意实值线性泛函, 且存在 $M$ 使得上述不等式成立, 则存在 $L^{q}(E)$ 中的唯一函数 $g$, 使得 $T$ 由积分给出. $L^{p}(E)$ 中的函数序列 $\{f_{n}\}$ 称为弱收敛于 $L^{p}(E)$ 中的函数 $f$, 若
\begin{align*}
\text{对所有 }g\in L^{q}(E),\quad\lim_{n\to\infty}\int_{E}f_{n}\cdot g=\int_{E}f\cdot g.
\end{align*}
我们用 Riesz 表示定理和 Helly 的一个定理证明, 对于 $1<p<\infty$, $L^{p}(E)$ 中的任何有界序列有一个弱收敛的子序列. 作为该结果的许多推论的一个例子, 我们证明某种凸泛函的最小值的存在性.

\subfile{Chapter-08/section-01.tex}

\subfile{Chapter-08/section-02.tex}

\subfile{Chapter-08/section-03.tex}

\subfile{Chapter-08/section-04.tex}

\end{document}